% Степень разработанности темы (только в диссертации); перенесено из бывшей главы 1
{\progress} Задача определения достаточного размера обучающей выборки и законы масштабирования изучаются в нескольких областях: в математической статистике, в байесовском анализе, в статистической теории обучения и в исследованиях ландшафта функции потерь нейронных сетей.

Особый интерес представляет вопрос о законах масштабирования параметрических моделей: по какому правилу следует совместно выбирать архитектуру модели и объем данных для ее обучения. В классических постановках задача определения достаточного размера выборки обычно рассматривается для линейных и обобщенно-линейных моделей, где модель задается конечномерным вектором параметров, а достаточность выборки формулируется через статистические критерии различимости гипотез о параметрах или через устойчивость вероятностной оценки модели~\cite{self1988,self1992,shieh2000,shieh2005,demidenko2007,motrenko2014}. Для обобщенно-линейных моделей используются, в частности, нормальная, логистическая и пуассоновская регрессии, объединяемые общей экспоненциальной формой распределения отклика и функцией связи между математическим ожиданием и линейным предиктором.

Классические статистические методы оценки размера выборки опираются на задание нулевой и альтернативной гипотез, допустимой вероятности ошибки первого рода и требуемой мощности теста. В наиболее общем виде достаточный размер выборки $m^*$ определяется как минимальное значение $m$, при котором статистика критерия $T_m$ обеспечивает требуемую мощность:
\[
    m^*:\quad \mathbb P(T_m > c_\alpha | H_1)\ge 1-\beta,
\]
где $c_\alpha$ "--- критическое значение уровня значимости $\alpha$, а $(1-\beta)$ "--- требуемая мощность. В этой постановке используются критерий отношения правдоподобия, статистика Вальда, метод множителей Лагранжа и их модификации~\cite{self1988,self1992,shieh2000,shieh2005,demidenko2007}. Эти методы дают теоретически обоснованные оценки, но требуют явных предположений о распределении данных, структуре модели и величине эффекта при альтернативе. На практике такие предположения часто недоступны или труднопроверяемы, особенно для многопараметрических нелинейных моделей.

Наряду со статистическими существуют байесовские и информационно-теоретические подходы. В байесовской постановке достаточность выборки связывается со стабилизацией апостериорного распределения параметров, уменьшением апостериорной неопределенности или малостью расстояния между распределениями, полученными по соседним подвыборкам. Типичными критериями являются средняя апостериорная дисперсия, среднее покрытие, средняя длина доверительной области и дивергенция Кульбака--Лейблера между апостериорными распределениями~\cite{lindley1997,rubin1998,wang2002,motrenko2014}. Информационно-теоретические подходы близки к этой линии и также используют энтропийные и дивергентные характеристики для описания связи между моделью и данными. Однако для моделей большой размерности такие методы часто оказываются вычислительно сложными и требуют дополнительных модельных предположений.

В статистической теории обучения размер выборки связывается со сложностью класса моделей через размерность Вапника"--~Червоненкиса (англ. Vapnik"--~Chervonenkis dimension, VC) и родственные меры емкости. Классические результаты этого типа лежат в основе анализа вероятно приблизительно корректного обучения (англ. probably approximately correct, PAC) и дают оценки числа объектов, достаточного для гарантированного обобщения. Однако для нейронных сетей VC-размерность, как правило, очень велика, вследствие чего получаемые оценки оказываются грубыми и существенно завышенными. Более современные меры, зависящие от выборки, такие как радемахеровские сложности, позволяют получать более тонкие оценки обобщающей способности. В русскоязычной литературе существенный вклад в развитие этой линии внес К.~В.~Воронцов, развивший комбинаторную теорию переобучения, методы получения точных оценок вероятности переобучения и оценки обобщающей способности, зависящие от выборки, в том числе в связи с радемахеровской сложностью и комбинаторными оценками надежности обучения по прецедентам~\cite{vorontsov2008pria,vorontsov2010dsc,vorontsov2012overfitting}. Для глубоких сетей были также предложены оценки через спектральную норму и отступ (англ. margin bounds), а также их уточнения, учитывающие произведения операторных норм весовых матриц~\cite{bartlett2017spectral}. Тем не менее даже такие оценки далеки от объяснения высокой обобщающей способности сильно перепараметризованных сетей, наблюдаемой на практике. Экспериментально показано, что современные глубокие сети способны точно аппроксимировать даже случайную разметку, что плохо согласуется с наивными комбинаторными оценками сложности~\cite{zhang2017rethinking}.

Классические подходы плохо переносятся на современные нейронные сети по нескольким причинам. Во-первых, комбинаторные и асимптотические меры сложности дают нереалистично большие или трудно интерпретируемые оценки необходимого объема данных. Во-вторых, явные статистические гипотезы о параметрах нейросетей обычно неизвестны. В-третьих, многократное полное переобучение модели на подвыборках разного размера для построения кривых обучения вычислительно дорого для больших сетей.

В то же время эмпирические исследования последних лет показывают существование устойчивых законов масштабирования: качество моделей систематически улучшается с ростом числа параметров, объема данных и вычислительного бюджета. В работе~\cite{kaplan2020scaling} показано, что для языковых моделей значение функции потерь хорошо аппроксимируется степенными законами по размеру модели, размеру выборки и объему вычислений. В работе~\cite{hoffmann2022} предложена постановка, оптимальная по вычислительному бюджету (англ. compute-optimal), согласно которой при фиксированном бюджете размер модели и число обучающих элементов последовательности (англ. tokens) должны масштабироваться согласованно; при этом многие крупные языковые модели оказались недообученными по данным. Эти результаты оказали значительное влияние на современное понимание компромисса между размером модели и объемом обучающей выборки, однако единого теоретического объяснения таких закономерностей для широкого класса параметрических моделей до сих пор нет.

Таким образом, возникает необходимость в развитии новых подходов к описанию достаточного размера выборки и законов масштабирования, опирающихся на свойства, доступные для анализа во всех параметрических моделях и инвариантные к выбору конкретного минимума. В настоящей диссертации в качестве такого объекта предлагается использовать локальную геометрию задачи минимизации эмпирической функции потерь, то есть свойства ландшафта функции потерь в окрестности точки минимума. Метод оценки масштабируемости, который использует эту локальную геометрию, а не глобальную картину поверхности, называется ландшафтным.

В задачах машинного обучения оптимум эмпирической функции потерь часто не единствен. В линейной регрессии при невырожденной матрице признаков существует единственный глобальный минимум, но при мультиколлинеарности или при $m < n$ задача вырождена: множество решений образует аффинное подпространство~\cite{seber2003,strang2006,vorontsov2010dsc}. В нейронных сетях функция потерь, как правило, невыпукла: существуют множественные локальные минимумы, седловые точки и плато~\cite{dauphin2014,choromanska2015,goodfellow2015qualitatively,li2018visualizing}. Более того, из-за симметрий (перестановка нейронов, знаковые преобразования) многие минимумы эквивалентны с точки зрения предсказательной способности~\cite{draxler2018essentially,entezari2022role,ainsworth2023git}. В таких условиях корректнее говорить не о единственной точке минимума, а о множестве оптимальных параметров модели.

Оптимизационной поверхностью или ландшафтом функции потерь (англ. loss landscape) называется график эмпирической функции потерь $\mathcal{L}_m(\mathbf{w})$, вычисленной по выборке размера $m$, как функции от параметров $\mathbf{w} \in \mathbb{R}^N$. Анализ этой поверхности для нейронных сетей является ключевой задачей теории оптимизации, что связано с итерационной природой процесса обучения "--- алгоритмы на основе стохастического градиентного спуска (англ. stochastic gradient descent, SGD) дают приближенное решение задачи минимизации эмпирической функции потерь, опираясь на глобальные и локальные свойства.

Исследование ландшафта функции потерь позволило выявить ряд важных характеристик, определивших развитие современных методов численной оптимизации~\cite{kingma2015adam,chaudhari2019,foret2021sam} "--- наличие множественных минимумов, седловых точек, многообразий меньшей размерности, содержащих множества оптимальных параметров~\cite{draxler2018essentially,garipov2018loss}, и др. В теории обобщения нейронных сетей плоские минимумы (с малой кривизной) часто связывают с лучшей обобщающей способностью~\cite{chaudhari2019,keskar2017largebatch,foret2021sam}, а узкие и глубокие "--- наоборот. Форма поверхности при этом определяется как самой архитектурой модели, так и заданной обучающей выборкой. Глубина и ширина нейронной сети давно перестали быть определяющими факторами для достижения высокого качества~\cite{he2016resnet}, из-за чего возникают новые архитектурные решения, позволяющие снизить сложность модели при неухудшении ее обобщающих свойств~\cite{ioffe2015batchnorm,ba2016layernorm,ulyanov2016instancenorm,wu2018groupnorm,srivastava2014dropout,henry2020qknorm}.

Ключевой характеристикой для анализа оптимизационной поверхности является матрица Гессе функции потерь $\mathbf{H}^{(m)}(\mathbf{w}) = \nabla^2_{\mathbf{w}} \mathcal{L}_{m}(\mathbf{w})$, которая содержит информацию о поведении ландшафта во втором порядке. Спектр матрицы Гессе характеризует нелинейность поверхности в окрестности минимума: большие собственные значения соответствуют  <<острым>> направлениям, малые "--- <<плоским>>~\cite{hochreiter1997flat,dauphin2014,keskar2017largebatch,li2018visualizing}. Более того, для типичных нейросетей наблюдается характерная структура спектра: множество малых собственных значений и небольшое число выбросов~\cite{papyan2019,ghorbani2019investigation,sankar2021deeper,singh2021analytic}. Множество работ посвящено еще более детальному анализу матрицы Гессе для различных архитектур моделей: выводу его явных выражений для полносвязных сетей~\cite{bishop1992hessian,pearlmutter1994,martens2012curvprop} и даже механизмов внимания в трансформерных сетях~\cite{ormaniec2024transformerhessian}.

Матрицы Гессе используются в методе Ньютона и в квазиньютоновских методах~\cite{nocedal2006numerical,dennis1983}. Наиболее известен среди последних алгоритм Бройдена"--~Флетчера"--~Гольдфарба"--~Шанно (англ. Broyden"--~Fletcher"--~Goldfarb"--~Shanno, BFGS)~\cite{broyden1970,fletcher1970}. Матрица Гессе и произведения Гессе на вектор применяются не только для построения методов второго порядка, но и для анализа локальной кривизны ландшафта, оценки <<остроты>> и <<плоскости>> минимумов, спектрального исследования динамики обучения, приближенного вычисления доверительных характеристик параметров, прореживания сети (англ. pruning) и малоранговых приближений (англ. low-rank), а также для построения гаусс"--~ньютоновских методов и методов без явной матрицы Гессе (англ. hessian-free) в больших нейросетевых моделях~\cite{martens2010hf,martens2012curvprop,papyan2019,zhang2024whytransformers}.

Несмотря на интерес научного сообщества к исследованию матрицы Гессе и ее применений в обучении нейронных сетей, вопрос о связи характеристик второго порядка оптимизационной поверхности с законами масштабирования остается открытым.
